\chapter{Urine Citrate}

\section*{What Is This Test?}
Urine citrate measures citrate excretion, usually in a timed urine collection. Citrate binds urinary calcium and inhibits crystal formation, so low citrate is an important risk factor for calcium kidney stones.

\section*{Alternative Names}
Urinary Citrate; Hypocitraturia Test; 24-Hour Urine Citrate.

\section*{Why This Test Is Ordered}
Testing is part of the metabolic evaluation of recurrent or high-risk kidney stones, nephrocalcinosis, and selected disorders causing chronic metabolic acidosis. It is commonly obtained with urine calcium, oxalate, uric acid, sodium, volume, and creatinine.

\section*{How This Test Is Performed}
The patient collects all urine over a defined interval, commonly 24 hours. The laboratory measures citrate concentration and calculates total daily excretion after checking collection completeness with urine volume and creatinine.

\section*{Preparation, Risks, and Limitations}
Follow collection and storage instructions carefully. The test is noninvasive. Diet, potassium depletion, metabolic acidosis, renal tubular disorders, chronic diarrhea, incomplete collection, and some medicines affect the result.

\section*{Understanding Results}
Low citrate indicates hypocitraturia and increased calcium-stone risk, although the laboratory threshold depends on sex, age, diet, and method. Results are interpreted with urine volume, calcium, oxalate, uric acid, sodium, pH, and stone composition.

\section*{References}
European Association of Urology guidance on metabolic evaluation of urolithiasis.

\clearpage
\section*{Understanding Results and Follow-up}
The laboratory report should identify the specimen, method, units, reference interval or decision limit, and whether the result is positive, negative, abnormal, or inconclusive. Reference intervals vary with age, sex, pregnancy, specimen type, assay, and laboratory; a result outside a range is not automatically a diagnosis, and a result within a range does not always exclude disease. Discuss unexpected results with the ordering clinician, who can decide whether repeat or confirmatory testing, treatment, or referral is appropriate. Do not change medicines or delay urgent care based on an isolated result.


\section*{Regulatory Status and Authorization Date}
Regulatory status applies to the specific assay, instrument, manufacturer, and intended use rather than to the test name alone. No single approval date can be assigned to this test category. For a commercial assay, verify the current product in the relevant regulator's database: the U.S. Food and Drug Administration (FDA) clearance, approval, or authorization; India's Central Drugs Standard Control Organisation (CDSCO) licence or permission; the European Union In Vitro Diagnostic Regulation (EU IVDR) CE certificate issued through the applicable conformity-assessment route; or the United Kingdom Medicines and Healthcare products Regulatory Agency (MHRA) registration and UKCA/CE status. Laboratory-developed tests may not have an individual FDA, CDSCO, EU IVDR, or MHRA approval date.
\clearpage
